\section{Finite-Dimensional Vector Spaces}

\subsection{Span and Linear Independence}

\begin{exercise}[2a1]
    Find a list of four distinct vectors in $\fbb^3$ whose span equals
    \[\set{(x, y, z) \in \fbb^3 \mid x + y + z = 0}\]
\end{exercise}

\begin{sol}
    Clearly, the span of the list involves two vectors, since we may rewrite $z$ in terms of $x$ and $y$ as $z = -x - y$.
	Thus, we may choose the following two vectors:
    \[v_1 = (1, 0, -1), \quad v_2 = (0, 1, -1).\]
    To find two more vectors, we may pick any two distinct linear combinations of $v_1$ and $v_2$.
	For example, we may select
    \[v_3 = v_1 + v_2 = (1, 1, -2), \quad v_4 = 2v_1 + 3v_2 = (2, 3, -5).\]
    Then
    \[\spn(v_1, v_2, v_3, v_4) = \spn(v_1, v_2) = \set{(x, y, z) \in \fbb^3 \mid x + y + z = 0}. \qh\]
\end{sol}

\begin{exercise}[2a2]
    Prove or give a counterexample: If $v_1, v_2, v_3, v_4$ spans $V$, then the list
    \[v_1 - v_2, v_2 - v_3, v_3 - v_4, v_4\]
    also spans $V$.
\end{exercise}

\begin{sol}
    Let $U = \set{v_1, v_2, v_3, v_4}$ and $W = \set{v_1 - v_2, v_2 - v_3, v_3 - v_4, v_4}$.
	Clearly, $v_4 \in \spn U$.
	Then we may write
    \begin{align*}
        v_3 &= (v_3 - v_4) + v_4 \in \spn W, \\
        v_2 &= (v_2 - v_3) + v_3 \in \spn W, \\
        v_1 &= (v_1 - v_2) + v_2 \in \spn W.
    \end{align*}
    Thus, $v_1, v_2, v_3, v_4 \in \spn W$, and so $V = \spn U \subseteq \spn W$.
	Since $W$ is a list of vectors in $V$, we have $\spn W \subseteq V$.
	Therefore, $\spn W = V$, and so $W$ spans $V$.
\end{sol}
 
\begin{exercise}[2a3]
    Suppose $v_1, \ldots, v_m$ is a list of vectors in $V$.
	For $k \in \set{1, \ldots, m}$, let
    \[w_k = v_1 + \cdots + v_k.\]
    Show that $\spn(v_1, \ldots, v_m) = \spn(w_1, \ldots, w_m)$.
\end{exercise}

\begin{sol}
    This is a generalization of \exref{2a2} and a direct consequence of the fact that $v_k = w_k - w_{k-1}$ for $k \in \set{2, \ldots, m}$ and $v_1 = w_1$.
	Since each $v_k$ is a linear combination of $w_1, \ldots, w_k$, we have $\spn(v_1, \ldots, v_m) \subseteq \spn(w_1, \ldots, w_m)$.
	On the other hand, each $w_k$ is a linear combination of $v_1, \ldots, v_k$, so $\spn(w_1, \ldots, w_m) \subseteq \spn(v_1, \ldots, v_m)$.
	Therefore, $\spn(v_1, \ldots, v_m) = \spn(w_1, \ldots, w_m)$.
\end{sol}
 
\begin{exercise}[2a4]
    \begin{subexercises}
        \item Show that a list of length one in a vector space is linearly independent if and only if the vector in the list is not 0.
        \item Show that a list of length two in a vector space is linearly independent if and only if neither of the two vectors in the list is a scalar multiple of the other.
    \end{subexercises}
\end{exercise}

\begin{sole}
    \item Let $v$ be the vector in the list.
	If $v = 0$, then the list is linearly dependent, since $0 \cdot v = 0$.
	If $v \neq 0$, then the only way to write $0$ as a linear combination of $v$ is to take the scalar multiple to be $0$.
	Thus, the list is linearly independent.
    \item Let $v_1$ and $v_2$ be the vectors in the list.
	If $v_1$ is a scalar multiple of $v_2$, then there exists a scalar $\lambda$ such that $v_1 = \lambda v_2$.
	Then we may write
    \[0 = v_1 - \lambda v_2,\]
    which shows that the list is linearly dependent.
	Conversely, if the list is linearly dependent, then there exist scalars $\alpha$ and $\beta$, not both zero, such that
    \[\alpha v_1 + \beta v_2 = 0.\]
    If $\alpha \neq 0$, then
    \[v_1 = -\frac{\beta}{\alpha} v_2.\]
    Similarly, if $\beta \neq 0$, then
    \[v_2 = -\frac{\alpha}{\beta} v_1.\]
    Thus, the list is linearly independent if and only if neither vector is a scalar multiple of the other.
\end{sole}
 
\begin{exercise}[2a5]
    Find a number $t$ such that
    \[(3, 1, 4), (2, -3, 5), (5, 9, t)\]
    is not linearly independent in $\mathbb{R}^3$.
\end{exercise}

\begin{sol}
    To say that the list is not linearly independent is to say that there exist scalars $\alpha, \beta, \gamma$, not all zero, such that
    \[\alpha(3, 1, 4) + \beta(2, -3, 5) + \gamma(5, 9, t) = (0, 0, 0).\]
    This is equivalent to the following system of equations:
    \begin{align*}
        0 & = 3\alpha + 2\beta + 5\gamma, \\
        0 & = \alpha - 3\beta + 9\gamma, \\
        0 & = 4\alpha + 5\beta + t\gamma.
    \end{align*}
    From the second equation, we obtain $\alpha = 3\beta - 9\gamma$.
	Substituting this into the first equation, we have
    \begin{align*}
        0 & = 3(3\beta - 9\gamma) + 2\beta + 5\gamma \\
        & = 11\beta - 22\gamma.
    \end{align*}
    Thus, $\beta = 2\gamma$.
	Substituting this into the expression for $\alpha$, we have $\alpha = 3(2\gamma) - 9\gamma = -3\gamma$.
	Substituting these values into the third equation, we have
    \begin{align*}
        0 & = 4(-3\gamma) + 5(2\gamma) + t\gamma \\
        & = -12\gamma + 10\gamma + t\gamma \\
        & = (t - 2)\gamma.
    \end{align*}
    Since $\gamma \neq 0$, we must have $t - 2 = 0$, or $t = 2$.
	Then the list is not linearly independent when $t = 2$.
\end{sol}
 
\begin{exercise}[2a6]
    Show that the list $(2, 3, 1), (1, -1, 2), (7, 3, c)$ is linearly dependent in $\mathbb{F}^3$ if and only if $c = 8$.
\end{exercise}

\begin{sol}
    We may proceed similarly to \exref{2a5} to obtain the following system of equations:
    \begin{align*}
        0 & = 2\alpha + \beta + 7\gamma, \\
        0 & = 3\alpha - \beta + 3\gamma, \\
        0 & = \alpha + 2\beta + c\gamma.
    \end{align*}
    The second equation gives us $\beta = 3\alpha + 3\gamma$.
	Substituting this into the first equation, we have
    \begin{align*}
        0 & = 2\alpha + (3\alpha + 3\gamma) + 7\gamma \\
        & = 5\alpha + 10\gamma.
    \end{align*}
    Thus, $\alpha = -2\gamma$.
	Substituting this into the expression for $\beta$, we have $\beta = 3(-2\gamma) + 3\gamma = -3\gamma$.
	Substituting these values into the third equation, we have
    \begin{align*}
        0 & = (-2\gamma) + 2(-3\gamma) + c\gamma \\
        & = -2\gamma - 6\gamma + c\gamma \\
        & = (c - 8)\gamma.
    \end{align*}
    Since $\gamma \neq 0$, we must have $c - 8 = 0$, or $c = 8$.
	Then the list is linearly dependent when $c = 8$.

    Conversely, if $c \neq 8$, then $(c - 8)\gamma = 0$ implies $\gamma = 0$.
	Then $\alpha = -2\gamma = 0$ and $\beta = -3\gamma = 0$, so the only solution to the system of equations is $\alpha = \beta = \gamma = 0$.
	Thus, the list is linearly independent when $c \neq 8$.
\end{sol}
 
\begin{exercise}[2a7]
    \begin{subexercises}
        \item Show that if we think of $\cbb$ as a vector space over $\rbb$, then the list $1 + i, 1 - i$ is linearly independent.
        \item Show that if we think of $\cbb$ as a vector space over $\cbb$, then the list $1 + i, 1 - i$ is linearly dependent.
    \end{subexercises}
\end{exercise}

\begin{sole}
    \item Suppose $\alpha, \beta \in \rbb$ such that
    \[\alpha(1 + i) + \beta(1 - i) = 0.\]
    Then we have
    \[(\alpha + \beta) + (\alpha - \beta)i = 0.\]
    Since $0 = 0 + 0i$, we must have $\alpha + \beta = 0$ and $\alpha - \beta = 0$.
	Solving this system of equations, we find that $\alpha = \beta = 0$.
	Thus, the list is linearly independent.
    \item Suppose $a + bi, c + di \in \cbb$ such that
    \[(a + bi)(1 + i) + (c + di)(1 - i) = 0.\]
    Then we have
    \[(a - b + c + d) + (a + b - c + d)i = 0.\]
    Then we have the system of equations
    \begin{align*}
        0 & = a - b + c + d, \\
        0 & = a + b - c + d.
    \end{align*}
    Adding the two equations results in $0 = 2a + 2d$, or $a = -d$.
	Substituting this into the first equation, we obtain $b = c$.
	Then we have infinitely many nontrivial solutions to the system of equations, such as $a = 1, b = 0, c = 0, d = -1$.
	Thus, the list is linearly dependent.
\end{sole}

\begin{exercise}[2a8]
    Suppose $v_1, v_2, v_3, v_4$ is linearly independent in $V$.
	Prove that the list
    \(v_1 - v_2, v_2 - v_3, v_3 - v_4, v_4\)
    is also linearly independent.
\end{exercise}

\begin{sol}
    Let $\alpha, \beta, \gamma, \delta \in \fbb$ such that
    \[\alpha(v_1 - v_2) + \beta(v_2 - v_3) + \gamma(v_3 - v_4) + \delta v_4 = 0.\]
    Then we have
    \[\alpha v_1 + (-\alpha + \beta)v_2 + (-\beta + \gamma)v_3 + (-\gamma + \delta)v_4 = 0.\]
    Since $v_1, v_2, v_3, v_4$ is linearly independent, we must have
    \begin{align*}
        0 & = \alpha, \\
        0 & = -\alpha + \beta, \\
        0 & = -\beta + \gamma, \\
        0 & = -\gamma + \delta.
    \end{align*}
    From the first equation, we have $\alpha = 0$.
	Substituting this into the second equation gives us $\beta = 0$.
	Substituting this into the third equation gives us $\gamma = 0$.
	Finally, substituting this into the fourth equation gives us $\delta = 0$.
	Thus, the only solution to the original equation is $\alpha = \beta = \gamma = \delta = 0$, and so the list is linearly independent.
\end{sol}
 
\begin{exercise}[2a9]
    Prove or give a counterexample: If $v_1, v_2, \ldots, v_m$ is a linearly independent list of vectors in $V$, then
    \(5v_1 - 4v_2, v_2, v_3, \ldots, v_m\)
    is linearly independent.
\end{exercise}

\begin{sol}
    Let $\alpha_1, \ldots, \alpha_m \in \fbb$ such that
    \[\alpha_1(5v_1 - 4v_2) + \alpha_2 v_2 + \cdots + \alpha_m v_m = 0.\]
    Then we have
    \[5\alpha_1 v_1 + (-4\alpha_1 + \alpha_2)v_2 + \alpha_3 v_3 + \cdots + \alpha_m v_m = 0.\]
    Since $v_1, v_2, \ldots, v_m$ is linearly independent, we must have
    \begin{align*}
        0 & = 5\alpha_1, \\
        0 & = -4\alpha_1 + \alpha_2, \\
        0 & = \alpha_3, \\
        & \vdots \\
        0 & = \alpha_m.
    \end{align*}
    From the first equation, we have $\alpha_1 = 0$.
	Substituting this into the second equation gives us $\alpha_2 = 0$.
	The remaining equations give us $\alpha_k = 0$ for $k = 3, \ldots, m$.
	Thus, the only solution to the original equation is $\alpha_k = 0$ for $k = 1, \ldots, m$, and so the list is linearly independent.
\end{sol}

\begin{exercise}[2a10]
    Prove or give a counterexample: If $v_1, v_2, \ldots, v_m$ is a linearly independent list of vectors in $V$ and $\lambda \in \fbb$ with $\lambda \neq 0$, then $\lambda v_1, \lambda v_2, \ldots, \lambda v_m$ is linearly independent.
\end{exercise}

\begin{sol}
    Let $\alpha_1, \ldots, \alpha_m \in \fbb$ such that
    \[\alpha_1(\lambda v_1) + \cdots + \alpha_m(\lambda v_m) = 0.\]
    Then we have
    \[\lambda(\alpha_1 v_1 + \cdots + \alpha_m v_m) = 0.\]
    Since $\lambda \neq 0$, we must have
    \[\alpha_1 v_1 + \cdots + \alpha_m v_m = 0.\]
    Since $v_1, v_2, \ldots, v_m$ is linearly independent, we must have $\alpha_k = 0$ for $k = 1, \ldots, m$.
	Thus, the only solution to the original equation is $\alpha_k = 0$ for $k = 1, \ldots, m$, and so the list is linearly independent.
\end{sol}
 
\begin{exercise}[2a11]
    Prove or give a counterexample: If $v_1, \ldots, v_m$ and $w_1, \ldots, w_m$ are linearly independent lists of vectors in $V$, then the list $v_1 + w_1, \ldots, v_m + w_m$ is linearly independent.
\end{exercise}

\begin{sol}
    Consider the following counterexample in $\rbb^2$:
    \[v_1 = (1, 0), \quad v_2 = (0, 1), \quad w_1 = (0, 1), \quad w_2 = (1, 0).\]
    Then $v_1, v_2$ and $w_1, w_2$ are linearly independent lists of vectors in $\rbb^2$.
	However,
    \[v_1 + w_1 = (1, 1), \quad v_2 + w_2 = (1, 1),\]
    and so the list $v_1 + w_1, v_2 + w_2$ is linearly dependent.
	Thus, the statement is false.
\end{sol}
 
\begin{exercise}[2a12]
    Suppose $v_1, \ldots, v_m$ is linearly independent in $V$ and $w \in V$.
	Prove that if $v_1 + w, \ldots, v_m + w$ is linearly dependent, then $w \in \spn(v_1, \ldots, v_m)$.
\end{exercise}

\begin{sol}
    Let $\alpha_1, \ldots, \alpha_m \in \fbb$, not all zero, such that
    \[\alpha_1(v_1 + w) + \cdots + \alpha_m(v_m + w) = 0.\]
    Then we have
    \[(\alpha_1 v_1 + \cdots + \alpha_m v_m) + (\alpha_1 + \cdots + \alpha_m)w = 0.\]
    If $\alpha_1 + \cdots + \alpha_m = 0$, then we have
    \[\alpha_1 v_1 + \cdots + \alpha_m v_m = 0,\]
    which contradicts the linear independence of $v_1, \ldots, v_m$.
	Thus, we must have $\alpha_1 + \cdots + \alpha_m \neq 0$.
	Let $\beta = \alpha_1 + \cdots + \alpha_m$.
	Then we may write
    \[w = -\frac{\alpha_1}{\beta}v_1 - \cdots - \frac{\alpha_m}{\beta}v_m.\]
    Thus, $w$ is a linear combination of $v_1, \ldots, v_m$, and so $w \in \spn(v_1, \ldots, v_m)$.
\end{sol}
 
\begin{exercise}[2a13]
    Suppose $v_1, \ldots, v_m$ is linearly independent in $V$ and $w \in V$.
	Show that
    \[v_1, \ldots, v_m, w \text{ is linearly independent} \iff w \notin \spn(v_1, \ldots, v_m).\]
\end{exercise}

\begin{sol}
    Suppose $v_1, \ldots, v_m, w$ is linearly independent.
	If there were scalars $\alpha_1, \ldots, \alpha_m, \beta$ such that
    \[\alpha_1 v_1 + \cdots + \alpha_m v_m + \beta w = 0,\]
    then linear independence would imply that $\alpha_1 = \cdots = \alpha_m = \beta = 0$.
	Suppose $w \in \spn(v_1, \ldots, v_m)$.
	Then there exist scalars $\gamma_1, \ldots, \gamma_m$ such that
    \[\gamma_1 v_1 + \cdots + \gamma_m v_m = w.\]
    Then we may write
    \[-\gamma_1 v_1 - \cdots - \gamma_m v_m + 1\cdot w = 0,\]
    which contradicts the linear independence of $v_1, \ldots, v_m, w$.
	Thus, $w \notin \spn(v_1, \ldots, v_m)$.

    Conversely, suppose $w \notin \spn(v_1, \ldots, v_m)$.
	Then there do not exist scalars $\alpha_1, \ldots, \alpha_m$ such that
    \[\alpha_1 v_1 + \cdots + \alpha_m v_m = w.\]
    If $v_1, \ldots, v_m, w$ is linearly dependent, then there exist scalars $\beta_1, \ldots, \beta_m, \gamma$, not all zero, such that
    \[\beta_1 v_1 + \cdots + \beta_m v_m + \gamma w = 0.\]
    If $\gamma = 0$, then we have
    \[\beta_1 v_1 + \cdots + \beta_m v_m = 0,\]
    which contradicts the linear independence of $v_1, \ldots, v_m$.
	Thus, we must have $\gamma \neq 0$.
	Then we may write
    \[w = -\frac{\beta_1}{\gamma}v_1 - \cdots -\frac{\beta_m}{\gamma}v_m,\]
    which contradicts the assumption that $w$ is not in the span of $v_1, \ldots, v_m$.
	Thus, $v_1, \ldots, v_m, w$ is linearly independent.
\end{sol}
 
\begin{exercise}[2a14]
    Suppose $v_1, \ldots, v_m$ is a list of vectors in $V$.
	For $k \in \set{1, \ldots, m}$, let
    \[w_k = v_1 + \cdots + v_k.\]
    Show that the list $v_1, \ldots, v_m$ is linearly independent if and only if the list $w_1, \ldots, w_m$ is linearly independent.
\end{exercise}

\begin{sol}
    Suppose $v_1, \ldots, v_m$ is linearly independent.
	Let $\alpha_1, \ldots, \alpha_m \in \fbb$ such that
    \[\alpha_1 w_1 + \cdots + \alpha_m w_m = 0.\]
    Substituting $w_k = v_1 + \cdots + v_k$ for each $k$, we get
    \[\alpha_1 v_1 + \alpha_2(v_1 + v_2) + \alpha_3(v_1 + v_2 + v_3) + \cdots + \alpha_m(v_1 + \cdots + v_m) = 0.\]
    Then we may rewrite this as
    \[(\alpha_1 + \cdots + \alpha_m) v_1 + (\alpha_2 + \cdots + \alpha_m) v_2 + \cdots + (\alpha_{m-1} + \alpha_m) v_{m-1} + \alpha_m v_m = 0.\]
    Since $v_1, \ldots, v_m$ is linearly independent, we must have $\alpha_1 = \cdots = \alpha_m = 0$.
	Hence, $w_1, \ldots, w_m$ is linearly independent.

    Conversely, suppose $w_1, \ldots, w_m$ is linearly independent.
	Let $\beta_1, \ldots, \beta_m \in \fbb$ such that
    \[\beta_1 v_1 + \cdots + \beta_m v_m = 0.\]
    Noting that $v_1 = w_1$ and $v_k = w_k - w_{k-1}$ for $k \in \set{2, \ldots, m}$, we may rewrite the equation as
    \begin{align*}
        \beta_1 w_1 + \beta_2 (w_2 - w_1) + \cdots + \beta_m (w_m - w_{m-1}) &= 0, \\
        (\beta_1 - \beta_2) w_1 + (\beta_2 - \beta_3) w_2 + \cdots + (\beta_{m-1} - \beta_m) w_{m-1} + \beta_m w_m &= 0.
    \end{align*}
    Since $w_1, \ldots, w_m$ is linearly independent, we find that the system of equations
    \begin{align*}
        0 & = \beta_1 - \beta_2, \\
        0 & = \beta_2 - \beta_3, \\
        & \vdots \\
        0 & = \beta_{m - 1} - \beta_m \\
        0 & = \beta_m,
    \end{align*}
    has the unique solution $\beta_1 = \cdots = \beta_m = 0$.
	Hence, $v_1, \ldots, v_m$ is linearly independent.
\end{sol}
 
\begin{exercise}[2a15]
    Explain why there does not exist a list of six polynomials that is linearly independent in $\pcal_4(\fbb)$.
\end{exercise}

\begin{sol}
    Since $\pcal_4(\fbb)$ can be spanned by the list $1, x, x^2, x^3, x^4$, and the length of a linearly independent list cannot exceed the length of a spanning list, there does not exist a list of six polynomials that is linearly independent in $\pcal_4(\fbb)$.
\end{sol}
 
\begin{exercise}[2a16]
    Explain why no list of four polynomials spans $\pcal_4(\fbb)$.
\end{exercise}

\begin{sol}
    Since the list $1, x, x^2, x^3, x^4$ is a linearly independent list in $\pcal_4(\fbb)$ and has length 5, and the length of a spanning list cannot be less than the length of a linearly independent list, no list of four polynomials can span $\pcal_4(\fbb)$.
\end{sol}
 
\begin{exercise}[2a17]
    Prove that $V$ is infinite-dimensional if and only if there is a sequence $v_1, v_2, \ldots$ of vectors in $V$ such that $v_1, \ldots, v_m$ is linearly independent for every positive integer $m$.
\end{exercise}

\begin{sol}
    Suppose $V$ is infinite-dimensional.
	We proceed by induction to construct a sequence of vectors $v_1, v_2, \ldots$ in $V$ such that $v_1, \ldots, v_m$ is linearly independent for every positive integer $m$.
	Since $V$ is infinite-dimensional, there exists a nonzero vector $v_1 \in V$.
	Suppose we have constructed a linearly independent list of vectors $v_1, \ldots, v_m$ in $V$.
	If every $v \in V$ is a linear combination of $v_1, \ldots, v_m$, then $V = \spn(v_1, \ldots, v_m)$, which contradicts the assumption that $V$ is infinite-dimensional.
	Thus, there exists a vector $v_{m+1} \in V$ that is not in $\spn(v_1, \ldots, v_m)$.
	Then by \exref{2a13}, the list $v_1, \ldots, v_m, v_{m+1}$ is linearly independent.
	By induction, we can construct a sequence of vectors $v_1, v_2, \ldots$ in $V$ such that $v_1, \ldots, v_m$ is linearly independent for every positive integer $m$.

    Conversely, suppose there exists a sequence of vectors $v_1, v_2, \ldots$ in $V$ such that $v_1, \ldots, v_m$ is linearly independent for every positive integer $m$.
	Then for every positive integer $m$, there exists a linearly independent list of vectors in $V$ of length $m$.
	Since there is no maximum length for a linearly independent list of vectors in $V$, this implies that $V$ is infinite-dimensional.
\end{sol}
 
\begin{exercise}[2a18]
    Prove that $\fbb^\infty$ is infinite-dimensional.
\end{exercise}

\begin{sol}
    By \exref{2a17}, it suffices to construct a sequence of vectors $v_1, v_2, \ldots$ in $\fbb^\infty$ such that $v_1, \ldots, v_m$ is linearly independent for every positive integer $m$.
	Let $v_k$ be the vector in $\fbb^\infty$ whose $k$-th coordinate is $1$ and all other coordinates are $0$.
	Explicitly, we have
    \[v_1 = (1, 0, 0, \ldots), \quad v_2 = (0, 1, 0, \ldots), \quad v_3 = (0, 0, 1, \ldots), \quad \ldots\]
    Then for any positive integer $m$, the list $v_1, \ldots, v_m$ is linearly independent, since the only way to write the zero vector as a linear combination of $v_1, \ldots, v_m$ is to take all coefficients to be zero.
	Thus, by \exref{2a17}, $\fbb^\infty$ is infinite-dimensional.
\end{sol}
 
\begin{exercise}[2a19]
    Prove that the real vector space of all continuous real-valued functions on the interval $[0, 1]$ is infinite-dimensional.
\end{exercise}

\begin{sol}
    By \exref{2a17}, it suffices to construct a sequence of continuous real-valued functions $f_1, f_2, \ldots$ on the interval $[0, 1]$ such that $f_1, \ldots, f_m$ is linearly independent for every positive integer $m$.
	Let $f_k(x) = x^{k-1}$ for each positive integer $k$.
	Then for any positive integer $m$, let $\alpha_1, \ldots, \alpha_m \in \rbb$ such that
    \[\alpha_1 f_1(x) + \cdots + \alpha_m f_m(x) = 0\]
    for all $x \in [0, 1]$.
	This is equivalent to the polynomial equation
    \[\alpha_1 + \alpha_2 x + \cdots + \alpha_m x^{m-1} = 0\]
    for all $x \in [0, 1]$.
	Since a polynomial of degree at most $m - 1$ can have at most $m - 1$ roots unless it is the zero polynomial, we must have $\alpha_1 = \cdots = \alpha_m = 0$.
	Thus, the list $f_1, \ldots, f_m$ is linearly independent.
	Therefore, by \exref{2a17}, the real vector space of all continuous real-valued functions on the interval $[0, 1]$ is infinite-dimensional.
\end{sol}
 
\begin{exercise}[2a20]
    Suppose $p_0, p_1, \ldots, p_m$ are polynomials in $\pcal_m(\fbb)$ such that $p_k(2) = 0$ for each $k \in \set{0, \ldots, m}$.
	Prove that $p_0, p_1, \ldots, p_m$ is not linearly independent in $\pcal_m(\fbb)$.
\end{exercise}

\begin{sol}
    Since $p_k(2) = 0$ for each $k \in \set{0, \ldots, m}$, we have
    \[p_k(x) = (x - 2)q_k(x)\]
    for some polynomial $q_k(x) \in \pcal_{m-1}(\fbb)$.
	Then we may write
    \[p_0(x), p_1(x), \ldots, p_m(x) = (x - 2)q_0(x), (x - 2)q_1(x), \ldots, (x - 2)q_m(x).\]
    Since $\pcal_{m - 1}(\fbb)$ can be spanned by the $m$ polynomials $1, x, \ldots, x^{m-1}$, the list of $m + 1$ polynomials $q_0, q_1, \ldots, q_m$ must be linearly dependent.
	Thus, there exist scalars $\alpha_0, \ldots, \alpha_m$, not all zero, such that
    \[\alpha_0 q_0(x) + \cdots + \alpha_m q_m(x) = 0.\]
    Then we may write
    \[\alpha_0 p_0(x) + \cdots + \alpha_m p_m(x) = (x - 2)(\alpha_0 q_0(x) + \cdots + \alpha_m q_m(x)) = 0.\]
    Since not all the scalars $\alpha_0, \ldots, \alpha_m$ are zero, this shows that $p_0, p_1, \ldots, p_m$ is linearly dependent in $\pcal_m(\fbb)$.
\end{sol}